\newcommand{\set}[1]{\left\{ \def\given{\ \middle| \ }  #1 \right\}  }
\newcommand{\defeq}{\mathrel{:=}}
\newcommand{\iso}{\cong}
\DeclareMathOperator{\id}{id}
\DeclareMathOperator{\tr}{tr}

\newcommand{\ii}{\mathsf{i}}

\newcommand{\transp}{\mathrm{T}}

%
%
\def\leftfun#1{\mathopen{}\left#1}
\def\rightfun#1{\right#1}

\NewDocumentCommand{\slg}{}{\operatorname{SL}_2(\mathbb{C})}
\NewDocumentCommand{\pslg}{}{\operatorname{PSL}_2(\mathbb{C})}
\NewDocumentCommand{\glg}{}{\operatorname{GL}_2(\mathbb{C})}
\NewDocumentCommand{\psp}{O{\mathbb{C}} O{1}}{{#1 P}^{#2}}

\NewDocumentCommand{\upf}{m}{\varphi^+\leftfun(#1\rightfun)}
\NewDocumentCommand{\dnf}{m}{\varphi^-\leftfun(#1\rightfun)}
\NewDocumentCommand{\aroundhol}{m}{\tau\leftfun(#1\rightfun)}

\NewDocumentCommand{\upr}{m}{#1^{\uparrow}}
\NewDocumentCommand{\dnr}{m}{#1^{\downarrow}}

\NewDocumentCommand{\gpoid}{m}{\Pi\leftfun(#1\rightfun)}

\NewDocumentCommand{\qndl}{}{\triangleleft}

\NewDocumentCommand{\Hol}{m}{\operatorname{Hol}_{#1}}
\NewDocumentCommand{\hol}{m}{\operatorname{hol}_{#1}}

\NewDocumentCommand{\evec}{m}{\mathrm{\mathbf{e}}_{#1}}

\DeclareMathOperator{\dil}{Li_2}
\DeclareMathOperator{\logdil}{L}


%
\ExplSyntaxOn
\DeclareExpandableDocumentCommand{\IfEmptyTF}{mmm}
 {
    \tl_if_empty:nTF {#1} {#2} {#3}
 }
\ExplSyntaxOff

\NewDocumentCommand{\potl}{o m}{
  \Phi_{\IfNoValueTF{#1}{}{#1}}%
  \IfEmptyTF{#2}{%
  }{%
    \leftfun(#2\rightfun)
  }
}
